
\fiche{Concordance du programme}{Repérer les contenus et les capacités attendues}
\phantomsection\label{concordance}\pdfbookmark[0]{Programme et sources}{programme}
Référence principale : enseignement optionnel de mathématiques complémentaires de Terminale générale, BO spécial n° 8 du 25 juillet 2019, arrêté du 19 juillet 2019, NOR MENE1921265A. Cette édition vise l'année 2026-2027.
Les intitulés ci-dessous sont des regroupements éditoriaux ; le document officiel reste la référence réglementaire.
\souscours{Analyse}
\noindent\textbf{A01} · Modéliser par suite explicite ou récurrente ; représentation en escalier. \hyperref[chap-1]{Chapitre 1}, exercices 5, 9, 11.\par\vspace{1.6mm}
\noindent\textbf{A02} · Limites finies et infinies, opérations, inégalités et gendarmes. \hyperref[chap-1]{Chapitre 1}, exercices 3, 8, 12.\par\vspace{1.6mm}
\noindent\textbf{A03} · Suites géométriques positives, sommes et limites de sommes. \hyperref[chap-1]{Chapitre 1}, exercices 2, 7.\par\vspace{1.6mm}
\noindent\textbf{A04} · Récurrence affine : équilibre, décalage et expression de toutes les solutions. \hyperref[chap-1]{Chapitre 1}, exercices 4, 5, 6.\par\vspace{1.6mm}
\noindent\textbf{A05} · Exponentielle, logarithme réciproque, limites et équation fonctionnelle. \hyperref[chap-2]{Chapitre 2}, exercices 1, 3, 7, 9.\par\vspace{1.6mm}
\noindent\textbf{A06} · Logarithme d’une puissance et recherche de seuils. \hyperref[chap-2]{Chapitre 2}, exercices 5, 6, 12.\par\vspace{1.6mm}
\noindent\textbf{A07} · Limites des fonctions usuelles ; asymptotes et formes indéterminées. \hyperref[chap-3]{Chapitre 3}, exercices 1, 2, 3, 10.\par\vspace{1.6mm}
\noindent\textbf{A08} · Continuité, valeurs intermédiaires, monotonie et nombre de solutions. \hyperref[chap-3]{Chapitre 3}, exercices 5, 6, 7, 9.\par\vspace{1.6mm}
\noindent\textbf{A09} · Réciproques : carré/racine et exponentielle/logarithme. \hyperref[chap-3]{Chapitre 3}, exercices 4.\par\vspace{1.6mm}
\noindent\textbf{A10} · Valeurs approchées, encadrement et résolution numérique. \hyperref[chap-3]{Chapitre 3}, exercices 8, 11, 12.\par\vspace{1.6mm}
\noindent\textbf{A11} · Dérivées : affine composée, carré de fonction, exponentielle et logarithme composés. \hyperref[chap-4]{Chapitre 4}, exercices 1, 2, 8.\par\vspace{1.6mm}
\noindent\textbf{A12} · Tableaux de variation, équations, inéquations, optimisation contextualisée. \hyperref[chap-4]{Chapitre 4}, exercices 4, 5, 6, 9, 11.\par\vspace{1.6mm}
\noindent\textbf{A13} · Convexité par sécantes et tangentes ; croissance de la dérivée, signe de la dérivée seconde. \hyperref[chap-5]{Chapitre 5}, exercices 1, 2, 6, 11.\par\vspace{1.6mm}
\noindent\textbf{A14} · Convexité, concavité, point d’inflexion et lecture de courbes. \hyperref[chap-5]{Chapitre 5}, exercices 3, 5, 9, 12.\par\vspace{1.6mm}
\noindent\textbf{A15} · Primitives, constante ; formes reconnues $2uu^{\prime}$, $u^{\prime}\mathrm e^u$, $u^{\prime}/u$. \hyperref[chap-6]{Chapitre 6}, exercices 1, 2, 3, 9.\par\vspace{1.6mm}
\noindent\textbf{A16} · Vérifier une solution ; résoudre $y^{\prime}=ay$ et $y^{\prime}=ay+b$. \hyperref[chap-6]{Chapitre 6}, exercices 4, 5, 6, 7, 12.\par\vspace{1.6mm}
\noindent\textbf{A17} · Intégrale positive ou signée, Chasles, primitive définie par intégrale. \hyperref[chap-7]{Chapitre 7}, exercices 1, 2, 4, 11.\par\vspace{1.6mm}
\noindent\textbf{A18} · Aires sous/entre courbes, encadrement, valeur moyenne et interprétation. \hyperref[chap-7]{Chapitre 7}, exercices 3, 5, 6, 12.\par\vspace{1.6mm}
\noindent\textbf{P01} · Conditionnement, probabilité totale, Bayes et répétitions conditionnelles. \hyperref[chap-8]{Chapitre 8}, exercices 2, 4, 5, 7, 10.\par\vspace{1.6mm}
\noindent\textbf{P02} · Uniforme discrète : loi, espérance ; Bernoulli : loi, espérance, écart type. \hyperref[chap-9]{Chapitre 9}, exercices 1, 2, 9.\par\vspace{1.6mm}
\noindent\textbf{P03} · Schéma de Bernoulli, arbre, coefficients, Pascal et symétrie. \hyperref[chap-9]{Chapitre 9}, exercices 3, 4, 5, 10.\par\vspace{1.6mm}
\noindent\textbf{P04} · Binomiale : masses, cumuls, espérance et écart type en résolution de problème. \hyperref[chap-9]{Chapitre 9}, exercices 5, 6, 7, 8, 11.\par\vspace{1.6mm}
\noindent\textbf{P05} · Intervalles binomiaux de grande ou faible probabilité ; décision et échantillonnage. \hyperref[chap-10]{Chapitre 10}, exercices 3, 5, 6, 11, 12.\par\vspace{1.6mm}
\noindent\textbf{P06} · Géométrique : loi, représentation, espérance, calcul et absence de mémoire. \hyperref[chap-11]{Chapitre 11}, exercices 1, 7, 11.\par\vspace{1.6mm}
\noindent\textbf{P07} · Densité, probabilité comme aire, fonction de répartition, intégrales d’espérance et de variance. \hyperref[chap-11]{Chapitre 11}, exercices 4, 5, 6, 12.\par\vspace{1.6mm}
\noindent\textbf{P08} · Uniforme continue $[0,1]$ puis $[a,b]$ : densité, répartition, espérance, variance. \hyperref[chap-11]{Chapitre 11}, exercices 2, 4, 9.\par\vspace{1.6mm}
\noindent\textbf{P09} · Exponentielle : densité, répartition, espérance et absence de mémoire. \hyperref[chap-11]{Chapitre 11}, exercices 3, 8, 10, 11.\par\vspace{1.6mm}
\noindent\textbf{S01} · Nuage et point moyen ; ajustement affine et moindres carrés. \hyperref[chap-12]{Chapitre 12}, exercices 1, 3, 5, 8, 12.\par\vspace{1.6mm}
\noindent\textbf{S02} · Corrélation, changement de variable, interpolation et extrapolation. \hyperref[chap-12]{Chapitre 12}, exercices 2, 4, 6, 7, 9, 10.\par\vspace{1.6mm}
\noindent\textbf{T01} · Fonctions Python, conditions, boucles, listes et programmation modulaire. \hyperref[chap-13]{Chapitre 13}, exercices 1, 2, 5, 6, 8, 11.\par\vspace{1.6mm}
\noindent\textbf{T02} · Ensembles, couples, connecteurs, négation, implication et réciproque. \hyperref[chap-13]{Chapitre 13}, exercices 3, 4, 7, 10.\par\vspace{1.6mm}
\noindent\textbf{T03} · Quantification, contre-exemple, raisonnement et communication orale. \hyperref[chap-13]{Chapitre 13}, exercices 9, 12.\par\vspace{1.6mm}

\fiche{Neuf thèmes, neuf dossiers}{Mettre les connaissances en situation}
\repere{projet-1}{Thème 1 · Modèles définis par une fonction : Une serre à construire}
\repere{projet-2}{Thème 2 · Modèles d’évolution : Deux modèles de refroidissement}
\repere{projet-3}{Thème 3 · Approche historique du logarithme : Du produit à la somme}
\repere{projet-4}{Thème 4 · Calculs d’aires : Une aire sans primitive usuelle}
\repere{projet-5}{Thème 5 · Répartition des richesses : Mesurer une inégalité}
\repere{projet-6}{Thème 6 · Inférence bayésienne : Une information qui change la décision}
\repere{projet-7}{Thème 7 · Répétitions et échantillonnage : Contrôler une surréservation}
\repere{projet-8}{Thème 8 · Temps d’attente : Attendre sans vieillir ?}
\repere{projet-9}{Thème 9 · Corrélation et causalité : Une croissance observée}

\souscours{Algorithmes proposés par le programme}
Termes récurrents et seuils : chapitres 1, 2 et 13. Balayage, dichotomie et Newton : chapitre 3.
Euler : chapitre 6. Rectangles, trapèzes et Monte-Carlo : chapitre 7 et dossier 4.
Coefficients binomiaux, cumuls et intervalles : chapitres 9 et 10. Simulation de Bernoulli, dé, géométrique, exponentielle et moyennes : chapitres 10 et 11, exercices associés et atelier ci-dessous.
Les exemples d'algorithmes et les pistes de problèmes du programme ne sont pas tous des exigences autonomes : la couverture des capacités est distinguée des prolongements.
\souscours{Démonstrations et histoire}
Somme géométrique : chapitre 1 ; identités et dérivée du logarithme : chapitre 2 ; primitives et équation différentielle : chapitre 6 ; lien aire/dérivée : chapitre 7 ; espérances uniforme et Bernoulli : chapitre 9 ; absence de mémoire : chapitre 11 ; moindres carrés : chapitre 12.
Les dossiers 3 et 4 relient les tables de logarithmes et les quadratures aux notions actuelles. Les travaux supplémentaires, notamment logistique ou Newton, sont guidés et distingués des connaissances exigibles.
\souscours{Évolution publiée pour septembre 2027}
Le nouveau programme de complémentaires, publié au BO n° 14 du 2 avril 2026, entre en application en 2027-2028. La comparaison des deux annexes retrouve les mêmes neuf thèmes et les mêmes grands blocs de contenus et capacités présentés ici. Le thème de corrélation propose aussi une concentration de médicament parmi les contextes possibles. Le préambule est réorganisé et le vocabulaire logique reste travaillé dans le manuel.
Cette édition est datée 2026-2027 : elle ne transforme pas les exemples possibles en obligations et n'anticipe pas des modalités d'évaluation non vérifiées.
\fiche{Atelier de simulation}{Compléter les algorithmes transversaux}
\begin{python}
from random import Random
from math import log, sqrt
rng = Random(2026)
def de():
    return 1 + int(6*rng.random())
def geometrique(p):
    k = 1
    while rng.random() >= p:
        k += 1
    return k
def exponentielle(taux):
    return -log(1-rng.random())/taux
def moyenne_des(n, repetitions):
    moyennes = [sum(de() for k in range(n))/n
                for j in range(repetitions)]
    m = sum(moyennes)/repetitions
    s = sqrt(sum((x-m)**2 for x in moyennes)/repetitions)
    seuil = 2*sqrt(35/12)/sqrt(n)
    proportion = sum(abs(x-3.5) <= seuil for x in moyennes)
    return m, s, proportion/repetitions
\end{python}
Utiliser $n\geq1$, un nombre de répétitions positif, $0<p\leq1$ et un taux strictement positif.
\code{de()} simule l'uniforme discrète à partir de l'uniforme continue. Pour une géométrique, on compte le premier succès ; pour une exponentielle, on inverse la fonction de répartition.
Comparer l'écart type simulé des moyennes à $\sqrt{35/12}/\sqrt n$ ; modifier le facteur $2$ en $3$ et commenter les proportions observées, sans les transformer en vérités exactes.
\fiche{Sources et repères éditoriaux}{Recherche vérifiée le 23 octobre 2026}
\begin{itemize}
\item Ministère, \href{https://www.education.gouv.fr/bo/19/Special8/MENE1921265A.htm}{arrêté du 19 juillet 2019 et annexe du programme de complémentaires}, BO spécial n° 8 du 25 juillet 2019.
\item Éduscol, \href{https://eduscol.education.gouv.fr/sites/default/files/document/spe265annexe1159134pdf-84162.pdf}{programme officiel en vigueur, 17 pages}.
\item Éduscol, \href{https://eduscol.education.gouv.fr/5817/programmes-et-ressources-en-mathematiques-voie-gt}{programmes et ressources en mathématiques, voie générale et technologique}, mise à jour juillet 2026.
\item Ministère, \href{https://www.education.gouv.fr/bo/2026/Hebdo14/MENE2902920A}{arrêté du 26 février 2026, BO n° 14 du 2 avril 2026}, entrée en application en 2027-2028 et annexe comparée.
\item Mission Villani-Torossian, \href{https://eduscol.education.gouv.fr/sites/default/files/document/grilles-de-positionnement-des-livres-de-mathematiques-70116.pdf}{grilles de positionnement des livres de mathématiques}, mai 2019 : structuration, calcul, compétences, modélisation et démonstrations.
\end{itemize}
Les cours et exercices sont rédigés pour ce manuel à partir des attendus. Les dossiers et les évaluations sont des créations originales : ils ne sont pas des sujets officiels d'examen. Les chiffres contextualisés sont fictifs, sauf identité mathématique explicitement démontrée.
